\documentclass{article}
\usepackage{graphicx} % Required for inserting images
\usepackage{hyperref}


\begin{document}

\begin{titlepage}
    \centering

    \vspace*{2cm}

    {\Huge\bfseries Portfolio Website and GitHub Repository Report\par}

    \vspace{2cm}

    {\Large B201 Computer Science Lab\par}

    \vspace{1cm}

    

    \begin{flushleft}
        \textbf{Student Name:} Eizeldeen Marzouq\\[0.5cm]
        \textbf{Student ID:} GH1060873\\[0.5cm]
        \textbf{Module:} B201 Computer Science Lab\\[0.5cm]
        \textbf{Lecturer:} William Baker Morrison\\[0.5cm]
        \textbf{University:} GISMA University of Applied Sciences\\[0.5cm]
        \textbf{Submission Date:} June 25 2026
    \end{flushleft}

    \vfill

\end{titlepage}



\tableofcontents
\newpage
\section{Introduction}
This report is for my final assessment for the module Computer science lab. It documents the steps taken to create the Cv which was done by latex, my personal portfolio website, which is for potential employers and recruiters to learn more about me, and lastly, my github repository and page that includes and hosts all my files.
\section{Portfolio structure}
The Portfolio website can be accessed here:
\href{https://eizeldeen-marzouq.github.io/EizeldeenMarzouq-Portfolio/#about}{Portfolio website}. The portfolio was designed as a single page webpage consisting of different subsections to make it easier to navigate and easier to understand. The subsections are the following:
\subsection{About}
This section gives the visitor a brief overview about who i am, what i do, and what i want to become. It is crucial to help the visitor understand me better.
\subsection{Education}
This provides the visitor with my previous studies along with my current studies at gisma university to ensure my qualifications.
\subsection{Experience}
This section highlights any previous work experience or projects i have done. This also allows the visitor to better understand my qualifications.
\subsection{Skills}
Next, the skills section highlight both my technical and soft skills. The technical skills include HTML and css while the soft skills include problem solving and teamwork.
\subsection{Languages}
This part highlights my proficiency in the 3 languages that i currently speak. A bar was included to make it better looking and easier to understand.
\subsection{Resume}
This sections contains a button that allows the visitor to view my Cv directly from the website.
\subsection{Contact}
This section also contains a button that allows the visitor to contact me by mail at anytime.
\section{Github Repository structure}
My Github Account can be accessed here:
\href{https://github.com/Eizeldeen-Marzouq}{Github Account}

\href{https://github.com/Eizeldeen-Marzouq/EizeldeenMarzouq-Portfolio}{Portfolio Repository}.
The Github repository stores all the files required for the website to work.
The repository includes the following main files and folders:
\begin{itemize}
    \item Index.html
    \item css
    \item Js
    \item img
    \item vendor
    \item cv
    
\end{itemize}
Index.html contains the main structure of the website, while the others store the designs and pictures used by the website.
Other files such as the README file were part of the original template and repository.
\section{Technologies used}
The following technologies were used in the making of the Cv, Portfolio, and the repository.
\subsection{HTML}
HTML was used to create the structure and the content of the website.
\subsection{CSS}
Css was used to customize the appearance of the website.
\subsection{Github}
Github was used to store and manage the the portfolio websites through a public repository.
\subsection{Github Pages}
Github pages was used to host the website and make it accessible online.
\subsection{Latex}
Latex was used to create the Cv and the following report.
\subsection{Other technologies}
Other technologies such as java script and Bootstrap were already used in the original template



\section{Design decisions}
Some design decisions were made during the making of the portfolio website.
Firstly,this specific template was chosen because i felt it was very easy to navigate and the visitor could access whatever information they were looking for quickly.
Next, language proficiency bars were added because it improves the visual appearance while also making it easier to understand. Next, a resume section included a view cv button instead of a download cv because it gives them the opportunity to view my cv without needing to download it, this decision was mainly due to the fact that most people don't prefer downloading random files. After that, a contact section was included with a direct email me button to ensure easy communication.
\section{Challenges encountered}
Many challenges were encountered during the making of the portfolio, cv, repository, and linking them all together. For example, sections like contact and resume didn't originally exist in the template, linking them to the navigation menu required some work but was fixed in the end. Another example would be when i deployed the website using github, the cv button wouldnt work due to difference in case sensitivity between the name of the file and the name of the reference in the code, after changing it, the issue was resolved. Other minor issues such as customizing buttons and bars were quickly resolved.
\section{Evaluation and Improvements}
\subsection{Strengths}
Overall, i think my portfolio, repository, and cv acheived its intended objective of presenting my professional and academic information to potential employers and recruiters. I believe my portfolio had the following strengths:
\begin{itemize}
    \item Professional layout
    \item easy to navigate
    \item easy access to cv
    \item easy access to contact information
\end{itemize}
\subsection{Improvements}
I also believe improvements are needed. Future improvements could include:
\begin{itemize}
    \item Adding more Software projects to my portfolio and github page
    \item including certifications
    \item adding more visuals instead of writing
    \item expanding my work experience section
\end{itemize}
\section{Conclusion}
The development of my portfolio and repository helped me gain important experience in website development, website deployment, github, and much more.
Through HTMl, Css, Github, Latex, and etc, a professional portfolio was created which showcases my experience, education, and skills which are crucial to potential employers and recruiters. The portfolio, Cv, and repository will continue improving in the future by adding more projects, certifications, experience, and education.
\end{document}
